% ============================================================
% Cheat sheet, argomento 10: package, componenti e deployment (dispensa, capitoli 8 e 9)
% ============================================================
\argomento{Package, componenti e deployment}{Dispensa: capitoli 8 e 9. Come è diviso il sistema e dove
gira}[Package, componenti, deployment]

\voce{package}{Package}{una cartella con la linguetta}
  {\begin{tikzpicture}
     \umlpkg{m}{(0,0)}{4.2}{1.3}{model}
     \node[classe1=1.3cm] at (1.1,0.45) {Borsa};
     \node[classe1=1.3cm] at (3.1,0.45) {Societa};
   \end{tikzpicture}}
  {Raggruppa classi (e altri package). In Java è un \code{package}: le classi della cartella \code{model} iniziano
   con \code{package model;}. Dentro si può disegnare il diagramma delle classi del package, o lasciarlo vuoto.}

\voce{dipendenzapackage}{Dipendenza fra package}{freccia tratteggiata: gli import}
  {\begin{tikzpicture}
     \umlpkg{g}{(0,0)}{1.6}{0.6}{gui}
     \umlpkg{c}{(2.3,0)}{1.6}{0.6}{controller}
     \umlpkg{m}{(4.6,0)}{1.4}{0.6}{model}
     \draw[dep] (g.east) -- (c.west);
     \draw[dep] (c.east) -- (m.west);
   \end{tikzpicture}}
  {\code{gui} dipende da \code{controller}, che dipende da \code{model}: nel codice sono gli \code{import}
   (\code{import model.Borsa;} in una classe del controller). Il \code{model} non importa nulla: si può cambiare
   l'interfaccia grafica senza toccarlo. Un \code{import} contrario alle frecce viola l'architettura.}

\voce{interfacciapackage}{Interfacce fra package}{realizzazione e dipendenza attraverso un'interfaccia}
  {\begin{tikzpicture}
     \umlpkg{d}{(0,0)}{2.5}{1.25}{dao}
     \node[classe1=2.1cm, align=center] (i) at (1.25,0.45) {{\tiny\normalfont«interface»}\\ListinoDAO};
     \umlpkg{p}{(3.2,0)}{2.6}{1.25}{postgres}
     \node[classe1=2.3cm, font=\tiny\bfseries] (k) at (4.5,0.45) {ListinoPostgresDAO};
     \draw[realizza] (k.west) -- (i.east);
   \end{tikzpicture}}
  {Il controller usa l'interfaccia \code{ListinoDAO} (package \code{dao}); la classe per PostgreSQL la realizza in un
   altro package. Per cambiare database si scrive un'altra implementazione, e chi usa \code{ListinoDAO} non cambia.}

\voce{componente}{Componente}{«component» e l'icona in alto a destra}
  {\begin{tikzpicture}
     \node[compo, minimum width=3.0cm, minimum height=1.2cm] (ts) {{\tiny «component»}\\ \textbf{Biglietteria}};
     \compicon{ts}
   \end{tikzpicture}}
  {Un pezzo del sistema che si costruisce e si sostituisce da solo: un eseguibile, una libreria (\code{.jar}), un
   servizio, un database. Contiene classi e package. Il diagramma dei componenti dice \textbf{che cosa} c'è e come si
   collega, \textbf{non dove} gira.}

\voce{interfaccecomp}{Interfacce fornite e richieste}{pallina, presa e il loro incastro}
  {\begin{tikzpicture}
     \node[compo, minimum width=1.6cm, minimum height=0.8cm] (c) at (0,0) {Vendite};
     \node[compo, minimum width=1.6cm, minimum height=0.8cm] (d) at (4.0,0) {Contabilità};
     \draw[line width=0.8pt] (c.east) -- (1.75,0);
     \presa{(1.95,0)}{0}
     \palla{(1.95,0)}
     \draw[line width=0.8pt] (2.06,0) -- (d.west);
     \node[ruolo] at (1.95,0.45) {Registro};
   \end{tikzpicture}}
  {\code{Contabilità} \textbf{fornisce} l'interfaccia \code{Registro} (pallina: la implementa); \code{Vendite} la
   \textbf{richiede} (presa: ne ha bisogno). L'incastro è il collegamento. In Java: una \code{interface}, una classe
   che la implementa e una che la riceve nel costruttore (\vv{componenticodice}).}

\voce{porta}{Porta}{quadratino sul bordo}
  {\begin{tikzpicture}
     \node[compo, minimum width=3.2cm, minimum height=1.3cm] (s) at (0,0) {};
     \compicon{s}
     \node[font=\tiny\bfseries, anchor=north] at ($(s.north)+(0,-0.08)$) {Server vendite};
     \node[compo, minimum width=1.4cm, minimum height=0.55cm, fill=white, font=\tiny] (t) at (0.3,-0.2) {Processore};
     \node[draw, fill=white, minimum size=0.16cm, inner sep=0pt] (p) at (-1.6,-0.2) {};
     \draw[dipend] (p.east) -- (t.west);
     \draw[line width=0.8pt] (p.west) -- (-2.2,-0.2);
     \palla{(-2.31,-0.2)}
   \end{tikzpicture}}
  {Il punto in cui i dati entrano o escono da un componente (con eventuali trasformazioni o serializzazioni). Un
   componente che ne contiene altri è un \textbf{sottosistema} («subsystem»): la porta inoltra all'interno ciò che
   arriva dall'esterno.}

\voce{nodo}{Nodo}{parallelepipedo: una macchina o un dispositivo}
  {\begin{tikzpicture}
     \nodo{n}{(0,0)}{3.6}{0.75}{\underline{sede : ServerBiglietti}}
   \end{tikzpicture}}
  {Un elemento \textbf{fisico}: un server, un PC, un telefono, un sensore. Il nome è un'istanza,
   \underline{\code{nome : Tipo}}. Un nodo può contenere \textbf{ambienti di esecuzione} (un DBMS, un container
   EJB, Docker), anche loro disegnati come nodi.}

\voce{artefatto}{Artefatto e componenti nel nodo}{«artifact»: un file prodotto o usato}
  {\begin{tikzpicture}
     \nodo{b}{(0,0)}{5.0}{1.7}{\textbf{Server della banca}}
     \node[compo, minimum width=2.2cm, minimum height=0.7cm, fill=white] (m) at (1.4,0.65) {{\tiny «component»}\\Mutui};
     \node[compo, minimum width=1.9cm, minimum height=0.7cm, draw=q9, fill=q9!6] (d) at (3.8,0.65) {{\tiny «artifact»}\\clienti.db};
     \draw[dipend] (m.east) -- (d.west);
   \end{tikzpicture}}
  {Dentro un nodo si disegnano i \textbf{componenti} che vi girano e gli \textbf{artefatti}: file eseguibili (\code{.jar},
   \code{.war}), database, file di configurazione. Così il diagramma di deployment si collega a quello dei
   componenti.}

\voce{connessione}{Connessione e valori etichettati}{una linea con il protocollo; \{nome = valore\}}
  {\begin{tikzpicture}
     \nodo{c}{(0,0)}{1.6}{0.9}{\textbf{PC}\\{\tiny browser}}
     \nodo{w}{(3.0,0)}{2.6}{0.9}{\textbf{Web server}\\{\tiny\{OS = Linux\}}}
     \draw[line width=0.8pt] (c.east) -- node[font=\tiny, above]{HTTPS} (w.west);
   \end{tikzpicture}}
  {Le connessioni sono collegamenti \textbf{fisici}, etichettati con il protocollo (TCP/IP, HTTP, JDBC, Java RMI). I
   \textbf{valori etichettati} \code{\{nome = valore\}} aggiungono dettagli: il sistema operativo, il software,
   quanti nodi uguali ci sono (\code{\{number deployed = 3\}}).}

\begin{caso}{componenticodice}{Interfaccia fornita e richiesta, nel codice}
\begin{lstlisting}[style=javacs]
interface Registro {          // l'interfaccia
    void registra(String vendita);
}
class Contabilita implements Registro {  // fornisce
    public void registra(String vendita) { }
}
class Vendite {
    private final Registro registro;     // richiede
    Vendite(Registro r) { registro = r; }
    void vendi(String articolo) {
        registro.registra(articolo);
    }
}
public class Assemblaggio {
    public static void main(String[] args) {
        Vendite v = new Vendite(new Contabilita());
    }
}
\end{lstlisting}
\tcblower
\begin{tikzpicture}
  \node[compo, minimum width=1.8cm, minimum height=0.9cm] (c) at (0,0) {Vendite};
  \compicon{c}
  \node[compo, minimum width=1.8cm, minimum height=0.9cm] (d) at (4.6,0) {Contabilita};
  \compicon{d}
  \draw[line width=0.8pt] (c.east) -- (2.1,0);
  \presa{(2.3,0)}{0}
  \palla{(2.3,0)}
  \draw[line width=0.8pt] (2.41,0) -- (d.west);
  \node[ruolo] at (2.3,0.45) {Registro};
\end{tikzpicture}
\regola{La \textbf{pallina} è la classe che \code{implements} l'interfaccia; la \textbf{presa} è la classe che la
riceve (nel costruttore) e la usa. L'incastro avviene dove si assemblano i pezzi (\code{main}): è l'unico punto che
conosce tutti e due. È la stessa idea della Dependency Injection (\vv{iniezione}).}
\end{caso}

\begin{attenzione}{Package, componenti o deployment?}
\begin{itemize}[nosep, leftmargin=*]
  \item \textbf{Package}: come è \textbf{organizzato il codice} (cartelle, \code{import}). Si sceglie con l'architettura,
        poi si disegna un diagramma delle classi per package.
  \item \textbf{Componenti}: \textbf{che cosa} è stato costruito e chi offre quale interfaccia a chi, in modo logico.
  \item \textbf{Deployment}: \textbf{dove} gira ogni pezzo, su quale macchina, con quale collegamento fisico. Lo stesso
        componente può comparire nel diagramma dei componenti e, dentro un nodo, in quello di deployment.
\end{itemize}
\end{attenzione}
